\documentclass[10pt,a4paper]{amsart}
\usepackage{amsmath,amssymb,amsthm,graphicx}
% For commutative diagrams we use Xy-pic
\usepackage[all]{xy}

% We use 2cell for 2-commutative diagrams.
\xyoption{2cell}
% Theorem environments.
%
\theoremstyle{plain}
\newtheorem{theorem}[subsection]{Theorem}
\newtheorem{proposition}[subsection]{Proposition}
\newtheorem{lemma}[subsection]{Lemma}

\theoremstyle{definition}
\newtheorem{definition}[subsection]{Definition}
\newtheorem{example}[subsection]{Example}
\newtheorem{exercise}[subsection]{Exercise}
\newtheorem{situation}[subsection]{Situation}

\theoremstyle{remark}
\newtheorem{remark}[subsection]{Remark}
\newtheorem{remarks}[subsection]{Remarks}

\numberwithin{equation}{subsection}

% Macros
%
\def\lim{\mathop{\mathrm{lim}}\nolimits}
\def\colim{\mathop{\mathrm{colim}}\nolimits}
\def\Spec{\mathop{\mathrm{Spec}}}
\def\Hom{\mathop{\mathrm{Hom}}\nolimits}
\def\Ext{\mathop{\mathrm{Ext}}\nolimits}
\def\SheafHom{\mathop{\mathcal{H}\!\mathit{om}}\nolimits}
\def\SheafExt{\mathop{\mathcal{E}\!\mathit{xt}}\nolimits}
\def\Sch{\mathit{Sch}}
\def\Mor{\mathop{\mathrm{Mor}}\nolimits}
\def\Ob{\mathop{\mathrm{Ob}}\nolimits}
\def\Sh{\mathop{\mathit{Sh}}\nolimits}
\def\NL{\mathop{N\!L}\nolimits}
\def\CH{\mathop{\mathrm{CH}}\nolimits}
\def\proetale{{pro\text{-}\acute{e}tale}}
\def\etale{{\acute{e}tale}}
\def\QCoh{\mathit{QCoh}}
\def\Ker{\mathop{\mathrm{Ker}}}
\def\Im{\mathop{\mathrm{Im}}}
\def\Coker{\mathop{\mathrm{Coker}}}
\def\Coim{\mathop{\mathrm{Coim}}}

% Boxtimes
%
\DeclareMathSymbol{\boxtimes}{\mathbin}{AMSa}{"02}

%
% Macros for moduli stacks/spaces
%
\def\QCohstack{\mathcal{QC}\!\mathit{oh}}
\def\Cohstack{\mathcal{C}\!\mathit{oh}}
\def\Spacesstack{\mathcal{S}\!\mathit{paces}}
\def\Quotfunctor{\mathrm{Quot}}
\def\Hilbfunctor{\mathrm{Hilb}}
\def\Curvesstack{\mathcal{C}\!\mathit{urves}}
\def\Polarizedstack{\mathcal{P}\!\mathit{olarized}}
\def\Complexesstack{\mathcal{C}\!\mathit{omplexes}}
% \Pic is the operator that assigns to X its picard group, usage \Pic(X)
% \Picardstack_{X/B} denotes the Picard stack of X over B
% \Picardfunctor_{X/B} denotes the Picard functor of X over B
\def\Pic{\mathop{\mathrm{Pic}}\nolimits}
\def\Picardstack{\mathcal{P}\!\mathit{ic}}
\def\Picardfunctor{\mathrm{Pic}}
\def\Deformationcategory{\mathcal{D}\!\mathit{ef}}
\usepackage{cleveref}
\crefname{originaldefinition}{original definition}{original definitions}
\Crefname{originaldefinition}{Original definition}{Original definitions}
\makeatletter
\newlabel{algebra-definition-ord}{{02MD}{source}{}{sourcepointer-1}{}}
\newlabel{algebra-definition-ord@cref}{{[originaldefinition][0][]02MD}{[1][0][]source}{}{}{}}
\makeatother
\title{Injectivity of a homomorphism of finite modules over a one-dimensional Noetherian local ring from the determinant at the prime, with the length of the cokernel given by an order}
\author{The Stacks Project Authors}
\date{stacks-project:chow.tex:local-snapshot-20261012; GNU FDL 1.2 or later}
\begin{document}
\maketitle
\noindent\textit{Editorial note.} All delivered mathematical text is the author's own, copied from the pinned Stacks Project chapter file chow.tex; no statement or proof step was rewritten, reordered or removed.\par\medskip
\noindent\textit{Editorial note.} Preamble provenance: the pinned chapter file opens with an input of the shared preamble and carries no preamble text of its own, so the project's single Stacks-wide preamble staged at sources/stacks-chow/preamble.tex is excerpted verbatim as two spans, lines 14-18 (Xy-pic setup) and lines 154-217 (the native theorem declarations and mathematical macros). Nothing was added, removed or redefined; the wrapper class and the amsmath/amssymb/amsthm/graphicx base packages are supplied by the assembler.\par\medskip
\noindent\textit{Editorial note.} The counted statement and proof contain no cross reference, no citation and no included file, so the delivered body is self contained.\par\medskip
\noindent\textit{Editorial note.} Canonical location of the counted proposition: Stacks tag 02QB (sources/stacks-chow/pinned-original-source.tex line 17804).\par\medskip
\noindent\textit{Allowed standard prerequisites (editorial).}\begin{itemize}
\item The order ord\_\{R/q\}(f) of an element in the residue field of a one-dimensional quotient. Algebra, Definition labelled algebra-definition-ord (Stacks tag 02MD), invoked at proof line 17838. Supplies the meaning of ord\_\{R/q\}(f) together with the identity length\_R(R/(f, q)) = ord\_\{R/q\}(f) used in the l = 1 case; the whole induction is carried out in the counted span.
\item The determinant det\_\{kappa(q)\}(N\_q) of a (2, 1)-periodic complex and its admissible relations. Same chapter, chow.tex, Section on periodic complexes and determinants, in particular the admissible relations (b) and (c) used at proof lines 17857, 17933 and 17945. Supplies the symbol notation and the relations that make the determinant comparison invariant under the changes of generating sequence performed in the proof.
\end{itemize}\medskip
\ifdefined\hypertarget\hypertarget{sourcepointer-1}{}\fi
\noindent\textit{Source reference.} Original definition 02MD: sources/stacks-chow/tags.txt, source lines 3440--3440. Cross-chapter standard prerequisite located in the project tag list staged in the local snapshot; invoked at proof line 17838 of the counted proof. It is retained only as a navigation pointer; every load-bearing step is inside the delivered span.\par\medskip
\begin{lemma}
\label{lemma-key-lemma}
Let $R$ be a Noetherian local ring.
Let $\mathfrak q \subset R$ be a prime with $\dim(R/\mathfrak q) = 1$.
Let $\varphi : M \to N$ be a homomorphism of finite $R$-modules.
Assume there exist $x_1, \ldots, x_l \in M$ and $y_1, \ldots, y_l \in M$
with the following properties
\begin{enumerate}
\item $M = \langle x_1, \ldots, x_l\rangle$,
\item $\langle x_1, \ldots, x_i\rangle / \langle x_1, \ldots, x_{i - 1}\rangle
\cong R/\mathfrak q$ for $i = 1, \ldots, l$,
\item $N = \langle y_1, \ldots, y_l\rangle$, and
\item $\langle y_1, \ldots, y_i\rangle / \langle y_1, \ldots, y_{i - 1}\rangle
\cong R/\mathfrak q$ for $i = 1, \ldots, l$.
\end{enumerate}
Then $\varphi$ is injective if and only if $\varphi_{\mathfrak q}$ is an
isomorphism, and in this case we have
$$
\text{length}_R(\Coker(\varphi)) = \text{ord}_{R/\mathfrak q}(f)
$$
where $f \in \kappa(\mathfrak q)$ is the element such that
$$
[\varphi(x_1), \ldots, \varphi(x_l)] = f [y_1, \ldots, y_l]
$$
in $\det_{\kappa(\mathfrak q)}(N_{\mathfrak q})$.
\end{lemma}

\begin{proof}
First, note that the lemma holds in case $l = 1$.
Namely, in this case $x_1$ is a basis of $M$ over $R/\mathfrak q$
and $y_1$ is a basis of $N$ over $R/\mathfrak q$ and we have
$\varphi(x_1) = fy_1$ for some $f \in R$. Thus $\varphi$ is injective
if and only if $f \not \in \mathfrak q$. Moreover,
$\Coker(\varphi) = R/(f, \mathfrak q)$ and hence the lemma
holds by definition of $\text{ord}_{R/q}(f)$
(see Algebra, Definition \ref{algebra-definition-ord}).

\medskip\noindent
In fact, suppose more generally that $\varphi(x_i) = f_iy_i$ for some
$f_i \in R$, $f_i \not \in \mathfrak q$. Then the induced maps
$$
\langle x_1, \ldots, x_i\rangle / \langle x_1, \ldots, x_{i - 1}\rangle
\longrightarrow
\langle y_1, \ldots, y_i\rangle / \langle y_1, \ldots, y_{i - 1}\rangle
$$
are all injective and have cokernels isomorphic to
$R/(f_i, \mathfrak q)$. Hence we see that
$$
\text{length}_R(\Coker(\varphi)) = \sum \text{ord}_{R/\mathfrak q}(f_i).
$$
On the other hand it is clear that
$$
[\varphi(x_1), \ldots, \varphi(x_l)] = f_1 \ldots f_l [y_1, \ldots, y_l]
$$
in this case from the admissible relation (b) for symbols.
Hence we see the result holds in this case also.

\medskip\noindent
We prove the general case by induction on $l$. Assume $l > 1$.
Let $i \in \{1, \ldots, l\}$ be minimal such that
$\varphi(x_1) \in \langle y_1, \ldots, y_i\rangle$.
We will argue by induction on $i$.
If $i = 1$, then we get a commutative diagram
$$
\xymatrix{
0 \ar[r] &
\langle x_1 \rangle \ar[r] \ar[d] &
\langle x_1, \ldots, x_l \rangle \ar[r] \ar[d] &
\langle x_1, \ldots, x_l \rangle / \langle x_1 \rangle \ar[r] \ar[d] &
0 \\
0 \ar[r] &
\langle y_1 \rangle \ar[r] &
\langle y_1, \ldots, y_l \rangle \ar[r] &
\langle y_1, \ldots, y_l \rangle / \langle y_1 \rangle \ar[r] &
0
}
$$
and the lemma follows from the snake lemma and induction on $l$.
Assume now that $i > 1$.
Write $\varphi(x_1) = a_1 y_1 + \ldots + a_{i - 1} y_{i - 1} + a y_i$
with $a_j, a \in R$ and $a \not \in \mathfrak q$ (since otherwise
$i$ was not minimal). Set
$$
x'_j =
\left\{
\begin{matrix}
x_j & \text{if} & j = 1 \\
ax_j & \text{if} & j \geq 2
\end{matrix}
\right.
\quad\text{and}\quad
y'_j =
\left\{
\begin{matrix}
y_j & \text{if} & j < i \\
ay_j & \text{if} & j \geq i
\end{matrix}
\right.
$$
Let $M' = \langle x'_1, \ldots, x'_l \rangle$ and
$N' = \langle y'_1, \ldots, y'_l \rangle$.
Since $\varphi(x'_1) = a_1 y'_1 + \ldots + a_{i - 1} y'_{i - 1} + y'_i$
by construction and since for $j > 1$ we have
$\varphi(x'_j) = a\varphi(x_i) \in \langle y'_1, \ldots, y'_l\rangle$
we get a commutative diagram of $R$-modules and maps
$$
\xymatrix{
M' \ar[d] \ar[r]_{\varphi'} & N' \ar[d] \\
M \ar[r]^\varphi & N
}
$$
By the result of the second paragraph of the proof we know
that $\text{length}_R(M/M') = (l - 1)\text{ord}_{R/\mathfrak q}(a)$
and similarly
$\text{length}_R(M/M') = (l - i + 1)\text{ord}_{R/\mathfrak q}(a)$.
By a diagram chase this implies that
$$
\text{length}_R(\Coker(\varphi')) =
\text{length}_R(\Coker(\varphi)) + i\ \text{ord}_{R/\mathfrak q}(a).
$$
On the other hand, it is clear that writing
$$
[\varphi(x_1), \ldots, \varphi(x_l)] = f [y_1, \ldots, y_l],
\quad
[\varphi'(x'_1), \ldots, \varphi(x'_l)] = f' [y'_1, \ldots, y'_l]
$$
we have $f' = a^if$. Hence it suffices to prove the lemma for the
case that $\varphi(x_1) = a_1y_1 + \ldots a_{i - 1}y_{i - 1} + y_i$,
i.e., in the case that $a = 1$. Next, recall that
$$
[y_1, \ldots, y_l] = [y_1, \ldots, y_{i - 1},
a_1y_1 + \ldots a_{i - 1}y_{i - 1} + y_i, y_{i + 1}, \ldots, y_l]
$$
by the admissible relations for symbols. The sequence
$y_1, \ldots, y_{i - 1},
a_1y_1 + \ldots + a_{i - 1}y_{i - 1} + y_i, y_{i + 1}, \ldots, y_l$
satisfies the conditions (3), (4) of the lemma also.
Hence, we may actually
assume that $\varphi(x_1) = y_i$. In this case, note that we have
$\mathfrak q x_1 = 0$ which implies also $\mathfrak q y_i = 0$.
We have
$$
[y_1, \ldots, y_l] =
- [y_1, \ldots, y_{i - 2}, y_i, y_{i - 1}, y_{i + 1}, \ldots, y_l]
$$
by the third of the admissible relations defining
$\det_{\kappa(\mathfrak q)}(N_{\mathfrak q})$. Hence we may
replace $y_1, \ldots, y_l$ by
the sequence
$y'_1, \ldots, y'_l =
y_1, \ldots, y_{i - 2}, y_i, y_{i - 1}, y_{i + 1}, \ldots, y_l$
(which also satisfies conditions (3) and (4) of the lemma).
Clearly this decreases the invariant $i$ by $1$ and we win by induction
on $i$.
\end{proof}

\end{document}
